\documentclass[11pt,a4paper]{article}
\usepackage[T1]{fontenc}
\usepackage{lmodern,amsmath,amssymb,amsthm,microtype,booktabs}
\usepackage[margin=27mm]{geometry}
\usepackage[colorlinks=true,linkcolor=blue,citecolor=blue,urlcolor=blue]{hyperref}
\setlength{\emergencystretch}{2em}
\newtheorem{theorem}{Theorem}[section]
\newtheorem{lemma}[theorem]{Lemma}
\newtheorem{proposition}[theorem]{Proposition}
\newtheorem{corollary}[theorem]{Corollary}
\newcommand{\Q}{\mathbb Q}\newcommand{\Z}{\mathbb Z}
\title{Restrictions on elliptic point-role couplings}
\author{}\date{17 November 2026}
\begin{document}\maketitle
\begin{abstract}
For the general Diophantine source, any fixed finite menu of nonzero
quadratic relations among its four root indices excludes infinitely
many positive inputs in every arithmetic progression. Fixed affine
couplings of the scalar and reference points also obstruct
completeness. Within the exact multiplication family, two independent
fixed group words have the same defect, including a finite menu of
pairs and all exceptional determinant fibers. Three individual
cross-root identifications are impossible. A shifted Lucas identity
supplies an integral single-coupling family with controllable fixed
finite congruences; its required joint prime and real realization
remains an additional arithmetic problem.
\end{abstract}

\section{Arithmetic setting}

Fix the arithmetic premises and a sufficiently large odd coefficient
exponent $L$ of the general source in \cite{General,Unit}. Put
\[
 K=\Q(\sqrt2),\quad D=11-6\sqrt2,\quad E:y^2=x^3-Dx,
 \qquad
 P=\left(\frac{57-28\sqrt2}{8},
          \frac{-660+417\sqrt2}{32}\right).
\]
The retained group statement is $2E(K)=\Z P$, with $P$ non-torsion.
The root and reference roles have nonzero integer labels
\[
 A=[n]P,\quad S=[a]P,\quad T=[t]P,\quad Z=[b]P,
 \qquad H=[h]P.
\]
Every guarded source solution satisfies
\begin{equation}\label{eq:source}
 \begin{gathered}
 n=\beta(X)=216X^4+450459,\qquad |a|>7,\qquad a\ne1,\qquad b\ne0,\\
 F_a(t,b):=t^2-a(a-2)tb+(a-2)^2b^2-1=0.
 \end{gathered}
\end{equation}
Thus $n$ is the input root index, whereas $a$ is the parameter of
the Lucas recurrence. The elementary Lucas lemma of \cite{Lucas}
gives
\begin{equation}\label{eq:growth}
                         \log|t|>|a|-3.
\end{equation}

Set
\[
 r_{\rm sc}(X)=\beta(\beta-1)(\beta-2401)(32\beta+343)(2\beta-1).
\]
There is a fixed finite exceptional set $\mathcal S_L$ such that,
for integral $X$ and $p\notin\mathcal S_L$ with
$v_p(r_{\rm sc}(X))=1$, scalar extraction gives
\begin{equation}\label{eq:scalar}
 \begin{gathered}
 p^3\mid a-1,\qquad H,\ S-P\text{ are formal at every place above }p,\\
                         v_p((a-1)/h)>0.
 \end{gathered}
\end{equation}
At the fixed prime $\ell=76249$, formality of both points and
$v_\ell((a-1)/h)>0$ also hold. This is a good split prime, and
$P$ has integral affine, hence nonidentity, reduction there.
These are source consequences; the unsquared comparison scale
preserves them by \cite{Unit}.

\begin{lemma}\label{lem:contacts}
Each of the five displayed factors of $r_{\rm sc}$ has arbitrarily
large prime divisors among its integer values. Outside a fixed
finite set, each such prime admits a full input class modulo its
square on which $v_p(r_{\rm sc}(X))=1$. On this class the
corresponding residue of $\beta(X)$ is, respectively,
\begin{equation}\label{eq:contacts}
                         0,\ 1,\ 2401,\ -343/32,\ 1/2.
\end{equation}
Any additional fixed finite set of primes may be avoided.
\end{lemma}
\begin{proof}
Schur's elementary argument supplies infinitely many prime divisors
of the values of any nonconstant integer polynomial. Indeed, after
an input translation its constant term $c$ is nonzero. If all prime
divisors lay in a finite set with product $N$, including the prime
divisors of $c$, its value at $cNu$, divided by $c$, would be an
integer congruent to one modulo $N$ and of absolute value greater
than one for sufficiently large $u$, a contradiction.

Here every factor is a nonzero constant plus a nonzero multiple of
$X^4$, and the five factors have disjoint roots. Their product is
squarefree over $\Q$. Outside its leading-coefficient and
discriminant primes, every residue root is simple. Exactly one of
its $p$ lifts modulo $p^2$ remains a root modulo $p^2$; any other
lift gives valuation one, while the other four factors remain units.
\end{proof}

\section{Quadratic restrictions on root indices}

\begin{lemma}\label{lem:bounded}
Let $R\in\Q[n,a,t,b]$ have total degree at most two.
There are fixed constants $C_0,C_1,C_2$ such that every integer
solution of \eqref{eq:source} and $R(n,a,t,b)=0$, with $b\ne0$,
$R(n,a,\cdot,\cdot)$ nonzero, and $|a|>C_2$, satisfies
\[
                       |a|\le C_0+C_1\log(2+|n|).
\]
\end{lemma}
\begin{proof}
Clear the fixed coefficient denominators of $R$.
For $|a|>7$, the homogenized conic $F_a$ is nonsingular, since
the determinant of its quadratic form is
\[
                   -(a-2)^2(1-a^2/4)\ne0.
\]
It is therefore geometrically irreducible. The coefficients $u,v$
of $t^2,tb$ in $R$ are fixed constants. If the nonzero specialized
polynomial $R(n,a,\cdot,\cdot)$ shares a component with $F_a$,
it equals $\lambda F_a$, and comparison gives
\[
                       \lambda=u,\qquad v=-u a(a-2).
\]
For $u=0$ this is impossible. For $u\ne0$ it permits only finitely
many $a$. Choose $C_2$ beyond their absolute values.

The resultant in $b$ is now a nonzero integer polynomial in $t$
vanishing at the given $t$. Use the actual specialized degree in
$b$ if that degree drops. Put $M=2+|a|+|n|$. All coefficient
polynomials of $F_a$ and $R$ in $t,b$ have height at most $CM^2$.
The Sylvester determinant has size at most four, so its coefficient
height is at most $C'M^8$. Its nonzero leading coefficient has
absolute value at least one. Cauchy's bound gives
$|t|\le1+C'M^8$. Comparing with \eqref{eq:growth} and absorbing
$O(\log(2+|a|))$ into $|a|/2$ proves the assertion.
\end{proof}

\begin{theorem}\label{thm:quadratic}
Fix a finite nonempty menu of nonzero polynomials in
$\Q[n,a,t,b]$ of total degree at most two. Require at least one
polynomial in the menu to vanish in the unchanged general source.
In every fixed arithmetic progression of inputs, infinitely many
positive integers admit no guarded rational witness. Finitely many
exceptional input repairs cannot restore completeness. The menu's
coefficients may depend on the fixed exponent $L$.
\end{theorem}
\begin{proof}
Start with a fixed input class $X\equiv X_*\pmod{M_*}$.
For each $R$, divide out its full power of $a-1$; this preserves
its vanishing at source labels. The resulting polynomial is
nonzero at $a=1$ as a polynomial in $n,t,b$. Some coefficient
$g(a,n)$ of its expansion in $t,b$ therefore satisfies
$g(1,n)\ne0$. Its degree in $n$ is at most two, so one of the
five contacts \eqref{eq:contacts} is not a root of $g(1,n)$.

Choose a distinct fresh prime for this relation at that contact,
using Lemma~\ref{lem:contacts}, and avoid the coefficient
denominators and the nonzero value of $g(1,n)$ there.
On its valuation-one input class, a source solution has
$a\equiv1$ and $n\equiv\beta(X)$ equal to the chosen contact.
Thus $g(a,n)$ is a unit, and $R(n,a,\cdot,\cdot)$ is nonzero.
Choose a prime separately for every menu member. The Chinese
remainder theorem combines these classes with the original class
into $X_0\pmod{M_0}$.

Choose arbitrarily large fresh scalar-contact primes $p\nmid M_0$.
Combine a valuation-one class modulo $p^2$ with the fixed class,
and take $0\le X<M_0p^2$. Every putative witness then satisfies
\[
 |a|\ge p^3-1,\qquad
 \log(2+|n|)=\log(2+\beta(X))\le C_3+8\log p.
\]
The first inequality excludes all exceptional $a$ from
Lemma~\ref{lem:bounded}, for large $p$. That lemma and the second
inequality give $|a|=O(1+\log p)$, a contradiction.
The constructed inputs include infinitely many distinct positive
integers: each fixed integral $X$ has $r_{\rm sc}(X)\ne0$ and
hence only finitely many prime divisors. All inputs lie in the
original progression.
\end{proof}

\section{The scalar point and the reference point}

\begin{proposition}\label{prop:reference}
For any fixed finite nonempty menu of integer pairs $(k,c)$,
there is an arithmetic progression of integer inputs on which
no guarded source witness satisfies any relation
\[
                             S=[k]H+[c]P
\]
from the menu. The progression may be chosen inside any
previously fixed input progression.
\end{proposition}
\begin{proof}
At a fresh scalar contact, $H$ and $S-P$ are formal.
For $c\ne1$, their indicated difference $[c-1]P$ must be
formal as well. Each of these finitely many fixed nonidentity
points is formal at only finitely many good primes. Exclude all
such primes. For $c=1$ and $k\ne0$, the index relation gives
\[
                     v_p((a-1)/h)=v_p(k).
\]
Exclude the prime divisors of these finitely many nonzero $k$;
this contradicts \eqref{eq:scalar}. The case $c=1,k=0$ gives
$a=1$, already excluded. A remaining fresh scalar prime yields
a full class modulo its square excluding every menu member.
Choose it prime to the previously prescribed modulus and combine
the two classes.
\end{proof}

\section{Restrictions on exact multiplication witnesses}

The exact multiplication family is
\begin{equation}\label{eq:prepared}
 U_A=[hn]P,\quad U_S=[ha]P,\quad
 U_T=[ht]P,\quad U_Z=[hb]P.
\end{equation}
These identities specify the witness family in this section.
They are not global conclusions for arbitrary source solutions.

\subsection{Individual cross-root identifications}

\begin{proposition}\label{prop:cross}
No guarded source witness in \eqref{eq:prepared} satisfies any of
\[
                          U_T=Z,\qquad U_Z=T,\qquad U_S=Z.
\]
\end{proposition}
\begin{proof}
The Pell relation gives $\gcd(t,b)=1$. If $ht=b$, then
$|t|=1$, contradicting \eqref{eq:growth}.
If $hb=t$, then $|b|=1$. The two roots $t$ and
$a(a-2)b-t$ of the resulting monic quadratic have positive
product $(a-2)^2-1$, so
\[
 |t|<|a(a-2)|\le |a|(|a|+2)<\exp(|a|-3)
 \qquad(|a|\ge8),
\]
again contradicting \eqref{eq:growth}.

Suppose $ha=b$. At one place over $\ell$, let $m>1$ be the
reduction order of the affine integral point $P$. Scalar
formality gives $m\mid h$, $m\mid a-1$, and hence $m\mid b$.
The elementary norm-one-unit descent of \cite{Lucas} writes
\[
                   |(a-2)b|=U_j(|a|),\qquad j\ge1,
\]
where $U_0(A)=0$, $U_1(A)=1$, and
$U_{j+1}(A)=AU_j(A)-U_{j-1}(A)$.
Since $|a|\equiv\pm1\pmod m$, this value is zero modulo $m$
precisely when $3\mid j$: at parameter $1$ the sequence is
$0,1,1,0,-1,-1$ periodically, and parameter $-1$ only changes
alternating signs. Thus $3\mid j$. The Lucas addition identity
gives $U_3(|a|)=a^2-1\mid U_j(|a|)$. Hence $a-1\mid b$,
because $\gcd(a-1,a-2)=1$. From $b=ah$ and
$\gcd(a,a-1)=1$, one obtains $a-1\mid h$, contradicting
$v_\ell((a-1)/h)>0$.
\end{proof}

\subsection{Finite menus of two group words}

A fixed group word is an integer linear combination of the nine
point roles and $P$. On \eqref{eq:prepared}, its vanishing becomes
\[
                         hB(n,a,t,b)+A(n,a,t,b)=0,
\]
where $A,B$ are affine forms with fixed rational coefficients.
The ten monomials
\[
                       1,n,a,t,b,h,hn,ha,ht,hb
\]
are linearly independent.
Two group words are consequently independent over $\Q$ exactly
when their pairs $(A_i,B_i)$ are independent over $\Q$.

\begin{theorem}\label{thm:two}
Fix a finite nonempty menu of pairs of fixed group-word equations
among the nine point roles and $P$, each pair linearly independent.
Require both equations from at least one pair to hold.
In every fixed arithmetic progression, infinitely many positive
integer inputs admit no guarded source witness in the exact
multiplication family satisfying this requirement. The conclusion
does not prescribe an order for choosing the root labels,
reference index, or comparison primes.
\end{theorem}
\begin{proof}
For a pair write its equations as $hB_i+A_i=0$ and put
$R=A_1B_2-A_2B_1$. When $R\ne0$, every witness satisfies a
nonzero quadratic relation among the root indices. Collect these
relations over the menu.

If $R=0$ and both $B_i$ vanish identically, collect a nonzero
affine form among the $A_i$. Otherwise reorder so that $B_1\ne0$.
If $A_1,B_1$ are not proportional over $\Q$, they are coprime
affine polynomials, with at least one of degree one. The identity
$A_1B_2=A_2B_1$ in the polynomial unique factorization domain
makes $A_2,B_2$ the same polynomial multiple of $A_1,B_1$.
The degree bound forces the multiplier to be constant, contrary
to independence of the pair. Therefore $A_1=cB_1$ for a rational
constant $c$, including $c=0$, and then $A_2=cB_2$. The equations
are $(h+c)B_i=0$.

Collect the nonzero affine form $B_1$ for this case. If $-c$ is
a nonzero integer, also list the possible fixed point $[-c]P$.
If $-c$ is not a nonzero integer, every witness already satisfies
$B_1=0$. This classifies every identically-zero determinant case.

The finite list of fixed nonidentity reference points is formal
at only finitely many good primes in total. Choose a fresh
scalar-contact prime outside their union and outside the modulus
of the prescribed input progression. On its valuation-one class,
\eqref{eq:scalar} excludes every fixed-reference possibility.
If the list is empty, no refinement is needed.
Every remaining witness must satisfy at least one of the collected
nonzero polynomials, all of total degree at most two. Apply
Theorem~\ref{thm:quadratic} on the refined input progression.
\end{proof}

Thus two point-role deletions by independent fixed group expressions
cannot preserve cofinite completeness within \eqref{eq:prepared},
even after changing the order of the construction or allowing a
fixed finite menu. The theorem does not concern arbitrary numerator
indices satisfying only local comparisons, nonconstant coefficients
in the group expressions, or non-group rational maps.

\section{A shifted Lucas coupling}

\begin{lemma}\label{lem:shifted}
For integers $a>2$ and $r\ge1$, define
\[
 j=6(a-2)r,\qquad
 b=\frac{U_j(a)}{a-2},\qquad
 t=U_{j+1}(a),\qquad
 h=\frac{U_j(a)}{(a-1)(a-2)}.
\]
Then $b,t,h$ are positive integers,
\begin{equation}\label{eq:shifted}
 F_a(t,b)=0,\qquad b=(a-1)h,\qquad
                         h\equiv-4r\pmod{a-1}.
\end{equation}
If $a$ is odd, then $b$ is even and $t$ is odd.
\end{lemma}
\begin{proof}
Since $U_j(a)\equiv j\pmod{a-2}$, the factor $a-2$ divides
$U_j(a)$. Since $3\mid j$, evaluation at parameter $1$ shows
that $a-1$ also divides $U_j(a)$. The two factors are coprime.
The recurrence is positive and increasing for $a>2$, and its
determinant identity gives
$U_{j+1}^2-aU_jU_{j+1}+U_j^2=1$.

For the stronger congruence, put
\[
 M(a)=\begin{pmatrix}a&-1\\1&0\end{pmatrix},\qquad
 J=\begin{pmatrix}1&-2\\2&-1\end{pmatrix}.
\]
Direct multiplication gives
\[
 M(a)^3\equiv-I+(a-1)J\pmod{(a-1)^2}.
\]
The lower-left entry of $M(a)^m$ is $U_m(a)$. Taking the $k$th
power consequently gives
\[
 \frac{U_{3k}(a)}{a-1}
              \equiv2k(-1)^{k+1}\pmod{a-1}.
\]
Here $k=2(a-2)r$ is even. Division by the unit $a-2$ modulo
$a-1$ proves \eqref{eq:shifted}. For odd $a$, the recurrence
modulo two has period $0,1,1$. Since $3\mid j$, $U_j(a)$ is
even and $U_{j+1}(a)$ is odd; division by the odd number $a-2$
preserves the parity of $b$.
\end{proof}

\begin{corollary}\label{cor:finite}
Let $M_0\ge1$ and $h_0\in\Z$ satisfy
$\gcd(4,M_0)\mid h_0$. Choose $r\ge1$ with
$-4r\equiv h_0\pmod{M_0}$. For every $a>2$ with
$a\equiv1\pmod{M_0}$, Lemma~\ref{lem:shifted} gives
$h\equiv h_0\pmod{M_0}$. At every prime $p$ such that
$v_p(a-1)>v_p(4r)$, one also has
\[
                  v_p(h)=v_p(4r),\qquad
                  v_p((a-1)/h)>0.
\]
\end{corollary}
\begin{proof}
Both conclusions follow directly from the congruence
$h\equiv-4r\pmod{a-1}$.
\end{proof}

The equation $b=(a-1)h$ realizes the single group coupling
$U_S=Z+H$ on the exact multiplication family.
Corollary~\ref{cor:finite} controls finite congruences prescribed
in advance and is compatible with positive scalar ratio valuations.
It does not assert the completeness of a coupled source.

Such completeness requires the actual Lucas roots and their trace
scales to be realized jointly with the comparison prime pattern.
At each selected prime the primary reduction order must divide
$h$, the opposite order must divide $h-1$, and the selected prime
must divide $h$, together with the required avoidance and
two-adic conditions. The finite and real scalar inequalities,
all denominator guards, and the pole preparation for the actual
trace inputs must hold simultaneously.
In the retained completeness proof, the roots and scales are fixed
before these primes, and $h$ is selected afterward. In
Lemma~\ref{lem:shifted}, $h$ is already fixed with $a,b,t$;
changing $r$ changes $b,t$ and the trace scales. The inherited
prime-pattern theorem \cite{Dependencies} and the fixed-modulus
congruence above do not supply this joint realization or the
required real returns of the coupled reference point.

\section{Exact algebra checks}

The companion \nolinkurl{check.py} verifies the conic determinant,
the universal quadratic-resultant identity and its coefficient
degrees for a general quadratic root relation, the matrix identity
in Lemma~\ref{lem:shifted}, and exact finite Lucas specializations.
It also checks the good split prime and affine reduction assertions
used in Proposition~\ref{prop:cross}. The checks use integer and
symbolic arithmetic and write a JSON certificate using
\nolinkurl{--output}. They certify these algebraic inputs, not the
retained arithmetic premises or the unresolved joint realization.

\begin{thebibliography}{9}
\bibitem{General}
\emph{A defining polynomial in twenty-five variables: combined scalar
and comparison conditions},
\nolinkurl{papers/few_variables25/few_variables25_merged_20261110_v19.tex}.
\bibitem{Unit}
\emph{A single norm determinant for the 25- and 13-variable constructions},
\nolinkurl{papers/finite_algebra/unit_norm_20261114_v28.tex}.
\bibitem{Lucas}
\emph{An explicit polynomial in 29 variables over the rationals},
\nolinkurl{papers/few_variables/few_variables_20261108_v8.tex},
the elementary Lucas growth and descent lemma.
\bibitem{Dependencies}
\emph{Arithmetic dependencies of the construction with 29 variables},
\nolinkurl{papers/arithmetic_dependencies/arithmetic_dependencies_20261108_v8.tex}.
\end{thebibliography}
\end{document}
